\begin{tabularx}{\textwidth}{Xrr}
    \rowcolor{heading} 1. Analysephase & & 4\  \\ \hline
    Ist-Analyse der bestehenden \ac{PVE}-Umgebung, Anforderungsaufnahme, Ressourcenplanung (Hardware, Storage) & 4 & \\ \hline
    \rowcolor{heading} 2. Planung \& Konzeption & & 5\  \\ \hline
    Erarbeitung des Backup-Konzepts (Strategie, Retention Policies, Netzwerk), Bewertung von Lösungsalternativen & 5 & \\ \hline
    \rowcolor{heading} 3. Beschaffung \& Vorbereitung & & 3\  \\ \hline
    Hardware bereitstellen, \ac{PBS}-ISO vorbereiten, \ac{VLAN} mit Netzwerkteam abstimmen und konfigurieren & 3 & \\ \hline
    \rowcolor{heading} 4. Installation \ac{PBS} & & 4\  \\ \hline
    \ac{PBS} auf dedizierter Hardware installieren, Grundkonfiguration (Netzwerk, Hostname, Updates) & 4 & \\ \hline
    \rowcolor{heading} 5. Konfiguration Datastore & & 4\  \\ \hline
    \ac{ZFS}-Pool und Backup-Datastore einrichten, Kapazitätsplanung und Quotas festlegen & 4 & \\ \hline
    \rowcolor{heading} 6. \ac{PVE}-Integration & & 4\  \\ \hline
    \ac{PBS} als Storage-Target in \ac{PVE} eintragen, Authentifizierung (\ac{API}-Token / Fingerprint) konfigurieren & 4 & \\ \hline
    \rowcolor{heading} 7. Backup-Jobs einrichten & & 5\  \\ \hline
    Backup-Jobs pro \ac{VM}-/\ac{CT}-Gruppe anlegen, Zeitpläne und Retention Policies definieren und testen & 5 & \\ \hline
    \rowcolor{heading} 8. Test \& Validierung & & 3\  \\ \hline
    Backup-Jobs manuell auslösen, Restore-Tests durchführen, Fehleranalyse und Nachbesserung & 3 & \\ \hline
    \rowcolor{heading} 9. Monitoring \& Benachrichtigung & & 2\  \\ \hline
    E-Mail-Benachrichtigungen konfigurieren, Monitoring-Integration prüfen & 2 & \\ \hline
    \rowcolor{heading} 10. Dokumentation & & 6\  \\ \hline
    Betriebsdokumentation erstellen: Installation, Konfiguration, Backup-Konzept, Restore-Verfahren & 6 & \\ \hline
    \rowcolor{heading} Summe & & 40 \\ \hline
\end{tabularx}
